% ===== §13 The Distribution of the Sensible =====
\section{The Distribution of the Sensible}
\label{sec:distribution}

\noindent
The third spine holds that the distribution of perceptual capacity is uneven in ways that constitute a neglected injustice, so that aesthetic education is answerable to the theory of justice. The paper has been laying this spine's foundations throughout: in the political economy of the senses, which showed perceptual capacity to be socially produced (\xref{sub:val-politicaleconomy}); in the account of value circulation, which distinguished the return of aesthetic value to its generators from its extraction (\xref{sub:val-circulation}); in the line between the refinement and the commodified training of perception (\xref{sub:plas-redline}); and in the production of perceptual space, which distributes the affordances of cultivation unequally (\xref{sub:space-home}). This section gathers those foundations into the claim of justice they were preparing. It sets out Ranci\`ere's account of the distribution of the sensible, which makes perception itself the primary political stake; it proposes, as the aesthetic instance of a form of injustice this series has developed, a generative aesthetic injustice; it confronts, as the paper's principal adversary in the register of political economy, Debord's spectacle; and it closes by returning to the claim with which the paper began, that beauty and ethics share a single root, and by gathering under a single criterion the ethics that the models of practice carried severally.

\subsection{Rancière and the distribution of the sensible}
\label{sub:dist-ranciere}

Ranci\`ere gave the political theory of perception its decisive formulation, and it is the frame this section requires. At the core of politics, he argued, there is an aesthetics, understood not as the annexation of politics by art but as the system of a priori forms determining what presents itself to sense experience: a distribution of the sensible that fixes, in advance of any particular political dispute, what is visible and what is not, what is audible and what goes unheard, who has a part in the common world and who is left without one \citep{ranciere2004}. The distribution of the sensible is at once a sharing, which establishes the perceptual conditions of a common world, and a partition, which divides those who are counted within it from those who are not perceived at all. Politics proper, for Ranci\`ere, occurs when those whom the reigning distribution renders imperceptible force a reconfiguration of it, demanding not merely to exist but to be perceived. This makes perception the primary stake rather than a secondary reflection of interests already formed elsewhere, and it converges exactly with the account this paper reached by another route. The distribution of the sensible is Ranci\`ere's name for the social determination of the perceptual frame that the convergence established (\xref{sec:convergence}): the a priori forms that decide what can be perceived are the frame, historically produced and unequally distributed, and Ranci\`ere adds that this distribution is itself the fundamental political fact. What a person can perceive, and whether a person is perceived, is distributed, and the distribution is a matter of justice. The paper's third spine is the aesthetics of the everyday read through this recognition: if perceptual capacity is distributed, and if the everyday is where the greater part of aesthetic life is generated, then the uneven distribution of the capacity to generate everyday aesthetic perception is a distribution of the sensible in Ranci\`ere's sense, and falls under justice as he showed perception to do.

\subsection{Generative aesthetic injustice}
\label{sub:dist-injustice}

The paper can now name the specific injustice its third spine concerns, and it does so by extending a form of injustice this series has already developed. Fricker identified epistemic injustice, the wrong done to a person specifically in their capacity as a knower, in two forms: the testimonial, in which prejudice deprives a speaker of due credibility, and the hermeneutical, in which a gap in shared interpretive resources leaves some experience unintelligible and inexpressible \citep{fricker2007}. This series has argued for a third, generative form of epistemic injustice, the wrong of being excluded from the generation of the interpretive resources themselves. The present paper proposes the aesthetic instance of this generative injustice. There is a generative aesthetic injustice, the wrong done to a person in their capacity as a generator of aesthetic perception, when they are deprived of the cultivation by which the everyday could be generated as perceived, or excluded from the shared generation of a perceived world, or reduced to the material of another's aesthetic experience without a share in the generating. This is not the injustice of being denied access to art, though it includes that; it is the deeper wrong of having one's perceptual life left uncultivated or extracted, one's capacity to generate the aesthetic perception of one's own ordinary days stunted by the social production of perceptual incapacity, or drawn off to sustain an aesthetic consumption in which one has no part. The political economy of the senses showed this injustice to be produced rather than natural: the person bound by need to a restricted sense, the person whose attention is captured and resold, the person whose home is built to starve perception and whose senses are trained for consumption, suffers a generative aesthetic injustice that a mode of life has inflicted. And the distinction between refinement and commodified training locates its mechanism: the same plasticity that could cultivate a perception whose value returns to the perceiver is turned, under the injustice, to training a perception whose value is extracted from them. Generative aesthetic injustice is the uneven and often deliberate distribution of the capacity to generate everyday aesthetic perception, and it is the specific wrong that the third spine holds aesthetic education answerable to.

\subsection{The spectacle}
\label{sub:dist-spectacle}

The paper's principal adversary in this register is the social form under which generative aesthetic injustice is now most pervasively produced, and it has a name and a theorist. Debord's spectacle is the condition in which social life, under modern conditions of production, presents itself as an immense accumulation of images, in which everything directly lived has receded into a representation, and in which the spectacle is, in his exact formulation, a social relation among people mediated by images \citep{debord1994}. Read through this paper's framework, the spectacle is the generalisation of the appropriative fate of the symbol (\xref{sub:sym-appropriative}) to the whole of social perception: it is the systematic substitution of the image for the perception, the outsourcing of seeing to a flow of representations, and the replacement of the generative encounter with the consumption of what is shown. The spectacle is thus the social engine of generative aesthetic injustice. It produces, at scale, the perceptual incapacity the third spine indicts: it habituates and numbs by supplying a perpetual stream of images that dispense with perceiving, it captures and resells the attention that everyday aesthetic life requires, and it trains perception for consumption under the guise of enriching it. Lefebvre's critique of everyday life supplies the register in which this operates, the everyday as the site where the reproduction of the relations of production is secured, and where consumer capitalism has colonised the ordinary that might otherwise have been the field of an unalienated life \citep{lefebvre1991}. The everyday is contested ground: it is at once the field of contingency in which generative aesthetic perception is possible and the territory that the spectacle colonises to produce a perception fitted to consumption. This is why the aesthetics of the everyday cannot be politically innocent. To cultivate the generative perception of ordinary life is, whether or not it intends to be, a refusal of the spectacle's substitution of the image for the seeing, and the paper's models of practice are, in this light, small counter-practices to a pervasive social production of perceptual incapacity. The paper does not claim that the arrangement of a shared home or the leaving of an emptiness overturns the spectacle; it claims that the cultivation of generative everyday perception is the practice through which, at the scale of a shared life, the spectacle's colonisation of the ordinary is locally refused.

\subsection{Beauty and ethics at one root}
\label{sub:dist-roots}

The paper returns, at its close, to the claim with which it began, that beauty and ethics share a single root, and it can now gather under one criterion the ethics its models carried severally. The fourth face of the paper's thesis held that aesthetic generation carries an ethical valence not by the later addition of a moral rule but at its root, because beauty is generated in a relation and a relation may return what it generates to all who stand in it or draw one of them off as fuel (\xref{sub:edu-dialectic}). The whole of the argument since has borne this out. The good and the bad symbolisation, the refinement and the commodified training, the good and the bad cycle of aesthetic value, the generative and the appropriative iteration of a symbolic strategy, and now the just distribution of the sensible and its spectacular corruption: these are one distinction drawn across the registers of the paper, and the single criterion beneath them is the one this series has carried throughout. An aesthetic generation is good, in perception as in symbol as in practice as in the distribution of perceptual capacity, when its value returns to all who generate it around an unsettled centre, a positive holonomy in which no one is reduced to the material of another's aesthetic life; and it is bad when its value is extracted, its centre settled, its cycle closed, its holonomy driven to zero or below, and some are made the unattended fuel of others' aesthetic consumption. The ethics that Model I carried in its sharing and its reproduction, that Model II built into its suspension of private judgement, that Model III drew from the dark virtue that generates without possessing, and that Model IV located in the direction of its refinement, are instances of this one criterion. And the justice that this spine concerns, the just distribution of the capacity to generate everyday aesthetic perception, is the same criterion at the scale of a society: a distribution is just when the capacity to generate the aesthetic perception of one's own ordinary life is not extracted, starved, or monopolised, but returned to all whose lives it is. Beauty and ethics are one at their root because the generation of the beautiful is a circulation of value, and the justice or injustice of a circulation, whether its surplus returns to its generators or is drawn off from them, is the ethical question as such. To cultivate the generative perception of the everyday is therefore already to take a side in a question of justice, and the aesthetic education this paper has proposed is, in the end, answerable to justice not as an external constraint upon it but as the truth of what it was all along.
